\chapter{Szoftver architektútra}

Ebben a fejezetben egy általános áttekintést adok a \textit{SuperSet} felépítéséről, a projekt struktúráról és a
fejlesztést támogató folyamatokról.

A \textit{SuperSet} egy elosztott tesz keretrendszser, amely Python 3.12 \cite{python312} nyelvben implementált. A
verzió kezelés Git segítségével történt. A \textit{SuperSet} mind logikailag mind modul szinten struktúrált. Ez a
struktúráltság teszi lehetővé a különböző absztrakciós szintek fenntartását és az implementációs részletek elrejtését.

\section{A projekt felépítés}

Ebben a fejezetben röviden ismertetem a szoftverkomponenseket amelyek felépítik a \textit{SuperSet} elosztott teszt
keretrendszert alkamazást. Mélyebb működésükre, konkrét implementációs részleteikre, és az azokat alátámasztó
indoklásokora a komponensek saját fejezeteiben fogok kitérni.

A teszt keretrendszer hét fő szoftver komponensből áll amelyek két \emph{nem} diszjunkt halmazba sorolhatók.

Az első halmaz a lokális futtatásért felelős szoftver komponenseket tartalmazza. Ebbe a halmazba tartozik a
\textit{run.py} fájl által implementált \textit{SuperSet} szoftver komponens, amely névadója lett a projektnek.

A második halmaz az elosztott futtatáshoz szükséges szoftver komponenseket tartalmazza. Ide tartozik a
\textit{client.py} fájl által implementált \textit{SuperSetConsumer} és a \textit{server.py} által implementált
\textit{SuperSetProducer} szoftver komponensek.

A két halmaz metszetében azok a szoftver komponensek szerepelnek amelyek elengedhetetlenek mind az elosztott, mind a
lokális futtatáshoz. Ezek a komponensek a \textit{Gekko}, a \textit{ConfigManager}, a \textit{SystemUnderTest}
(továbbiakban \textit{Sut}), illetve a program kimenet szabványos formára hozásáért felelős \textit{Reporter}.

\subsection{A \textit{SuperSet} szoftver komponens}

A \textit{SuperSet} szoftver komponens a lokális futtatásért felelős fő komponens. Feladata a \textit{ConfigManager}
szoftver komponens által előállított teszt csomagok futtatása, valamint a \textit{Gekko} szoftver komponens által
előállított \textit{fuzz} típusú teszt csomagok futtatása is. Továbbá a \textit{ConfigManagertől} megkapja a tesztekre
érvényes beállításokat is, azok kapcsolatatival és kizárásaival, így megalkotja a futtatási mátrixot. Kimenetét egy
\textit{SQLite} \cite{sqlite} adatbázisban tárolja.
